
Рассмотрим как меняется матрица билинейной формы при смене базиса. Пусть матрица $B$ задана в базисе $e$. Перейдем в базис $e'$ по средствам матрицы перехода $S$, тогда:
\[
x=Sx', \ \ y=Sy'
\]
Подставляя данные равенства в $b(x,y)=x^TBy$, получаем
\[
b(x,y)=x^TBy=(Sx')^TB(Sy')=(x')^TS^TBSy'
\]
откуда в силу единственности $B'$ следует $B'=S^TBS$. Из не вырожденности $S$ следует \[
rg(B')=rg(B), \ \ \det B'=\det B
\]

\begin{definition}
    Билинейная форма $b(x,y)$ называется симметричной, если \[
    \forall x,y\in L \hookrightarrow b(x,y)=b(y,x)
    \]
\end{definition}
\begin{proposition}
    Билинейная форма $b(x,y)$ является симметричной тогда и только тогда, когда соответствующая матрица $B$ является симметричной.
\end{proposition}
\begin{proof}  
    Пусть матрица $B$ задана в базисе  $e$.
    
    $\Rightarrow$ Из определения симметричности $\forall e_i, e_j\in e \hookrightarrow b(e_i,e_j)=b(e_j,e_i)$, что по определению означает симметричность матрицы $B$.

    $\Leftarrow$ Рассмотрим $b(x,y)$ в точке $(x,y)$ как матрицу размером $1\times 1$. Тогда $b(x,y)=b^T(x,y)$, следовательно $x^TBy=(x^TBy)^T=y^TB^Tx=b(y,x)$. Значит $b(x,y)=b(y,x)$.
\end{proof}
\subsection{Квадратичные формы}
\begin{definition}
    Пусть $b(x,y)$ -- симметричная билинейная форма. Тогда квадратичной формой называется \[
    k(x)=b(x,x)
    \]
\end{definition}
Рассмотрим вопрос о том, как выразить $b(x,y)$ через соответствующую квадратичную форму (в частности отсюда будет следовать единственность порождающей билинейной формы для квадратичной формы):
\[
k(x+y)=b(x,x)+b(x,y)+b(y,x)+b(y,y)=k(x)+2b(x,y)+k(y)
\]
\[
b(x,y)=\frac{k(x+y)-k(x)-k(y)}{2}
\]
\begin{definition}
    Матрицей квадратичной формы называется матрица соответствующей билинейной формы.
\end{definition}
\begin{example}
    Квадратичная форма $k(x)=x_1^2+3x_1x_2-7x_2^2$. Ее порождает билинейная форма $b(x,y)=x_1y_1+1.5x_1y_2+1.5x_2y_1-7x_2y_2$. Так же $k(x)$ можно представить в виде
    \[
    k(x)=\begin{pmatrix}
        x_1&x_2
    \end{pmatrix}\begin{pmatrix}
        1&1,5\\
        1.5&-7
    \end{pmatrix}\begin{pmatrix}
        x_1\\
        x_2
    \end{pmatrix}
    \]
\end{example}
\begin{definition}
    Квадратичная форма в базисе $e^{(n)}$ имеет диагональный вид, если \[
    k(x)=\sum_{i=1}^{n} \alpha_i x_i^2
    \]
\textit{Замечание}: в данном базисе матрица квадратичной формы диагональная.
\end{definition}
\begin{definition}
    Квадратичная форма в базисе $e^{(n)}$ имеет канонический вид, если \[
    k(x)=\sum_{i=1}^{n} \alpha_i x_i^2, \ \ \alpha_i\in \{0, -1, 1\}
    \]
\end{definition}
\begin{proposition}
    Из приводимости квадратичной формы к диагональному виду следует ее приводимость к каноническому виду.
\end{proposition}
\begin{proof}
    Если существует базис, в котором $k(x)=\sum_{i=1}^{n} \alpha_i x_i^2$, то делая замену $x_i'=\sqrt{|\alpha_i|}x_i$, получаем $ k(x)=\sum_{i=1}^{n} (x_i')^2sign(\alpha_i)$ -- канонический вид.
\end{proof}
Теперь покажем, что любую квадратичную форму можно привести к диагональному виду.
\begin{theorem}
    Для любой квадратичной формы существует базис, в котором она имеет диагональный вид.
\end{theorem}
\begin{proof}
    Пусть задана квадратичная форма $k(x)$ и её матрица $B$ в некотором базисе. Рассмотрим смену базиса при котором матрица перехода $S$ является элементарной. При таком переходе $B'=S^TBS$, т.е. сначала применяется (не умаляя общности) элементарная операция к столбцам матрица $B$ (соответствующая элементарной матрице  $S$), а потом к строкам $BS$ (соответсвующая матрице $S^T$).

    Применим к $B$ последовательность таких элементарных преобразований. Опишем алгоритм:
    \begin{enumerate}
        \item "Основной случай": если элемент матрицы $b_{11}\neq 0$, то с помощью метода Гаусса обнуляем все элементы в первом столбце. В силу сохранения симметрии $B$ при смене базиса обнуляться и все элементы первой строки. Получаем
        \[
        \begin{pmatrix}
            b_{11}&0&\dots&0\\
            0 &b_{22} &\dots &b_{2n}\\
            \vdots &\vdots & \ddots&\\
            0&b_{n2}&&b_{nn}
        \end{pmatrix}
        \]
        Далее переходим к рассмотрению подматрицы, получающейся исключением первой строки и первого столбца.
        \item "Особый случай": если элемент матрицы $b_{11}= 0$.
            \begin{itemize}
                \item Если первый столбец и первая строка целиком состоят из нулей, то переходим к рассмотрению подматрицы, получающейся исключением первой строки и первого столбца.
                \item Если же существует элемент $b_{k1}\neq 0$, то смотрим на элемент $b_{kk}$. 
                
                Если $b_{kk}\neq0$, то применяем преобразование переставляющее $1$ строку с $k$ и $1$ столбец с $k$, в результате чего $b_{kk}\to b_{11}$ и получаем основной случай.
                
                Если $b_{kk}=0$, то к $1$ строке прибавляем $k$ строку и получаем снова основной случай.
            \end{itemize}
    \end{enumerate}
    Таким образом, после $n$-ого шага получаем \[
    \begin{pmatrix}
            b_{11}&0&\dots&0\\
            0 &b_{22} &\dots &0\\
            \vdots &\vdots & \ddots&\\
            0&0&&b_{nn}
        \end{pmatrix}
    \]
\end{proof}
\textit{Следствие}: Для любой квадратичной формы существует базис, в котором она имеет канонический вид.

На практике для приведения к каноническому виду используется \hyperref[Lagranj]{метод Лагранжа}.
\begin{definition}
    Рангом квадратичной формы называется количество ненулевых элементов на диагонали в каноническом виде.
\end{definition}
Ранее мы показали, что $rg(B)$ не зависит от базиса. Значит ранг квадратичной формы совпадает с рангом соответствующей матрицы (т.е. $rg(k(x))=rg(B)$).

\begin{definition}
    Квадратичная форма $k(x)$ называется 
    \begin{itemize}
        \item положительно определенной на $L$, если 
        \[
        \forall x\neq0 \in L \hookrightarrow k(x)>0
        \]
        \item положительно полуопределенной на $L$, если 
        \[\forall x \in L \hookrightarrow k(x)\geqslant 0\]
        \item отрицательно определенной на $L$, если 
        \[
        \forall x\neq0 \in L \hookrightarrow k(x)<0
        \]
        \item отрицательно полуопределенной на $L$, если 
        \[\forall x \in L \hookrightarrow k(x)\leqslant 0\]
    \end{itemize}
\end{definition}
\textit{Замечание}: форма $k(x)=0$ и положительно полуопределена и отрицательно полуопределена.

\begin{definition}
    Пусть $L^{(-)}\subset L$ -- подпространство максимальной размерности, на котором $k(x)$ отрицательно определена. Тогда $\dim L^{(-)}$ называется отрицательным индексом квадратичной формы.
    
    Аналогично определяется положительный индекс квадратичной формы.
\end{definition}
\begin{theorem} (Закон инерции квадратичных форм)

    Количество положительных и отрицательных квадратов в каноническом виде квадратичной формы является инвариантом (относительно смены базиса).
\end{theorem}
\begin{proof}
    Приведем квадратичную форму к каноническому виду в базисе $f^{(n)}$
    \[
    k(x)=-(x_1')^2-\dots-(x_m')^2+(x_{m+1}')^2+\dots+(x_r')^2
    \]
    где $r=rg(k(x))$. Пусть $L_1=\langle f_1, \dots, f_m \rangle$, $L_2=\langle f_{m+1}, \dots, f_n \rangle$. Тогда $k(x)$ отрицательно определена на $L_1$ и положительно полуопределена на $L_2$, следовательно \[
    \forall x\neq 0\in L_1 \hookrightarrow k(x)<0 \quad  \forall x \in L_2 \hookrightarrow k(x)\geqslant 0
    \]
    Предположим, что $\dim L_1 <\dim L^{(-)}$. Тогда $\dim L^{(-)}\cap L_2 > m+\dim L_2=n$. Значит $L^{(-)}\cap L_2$ непустое, что противоречит определению $L^{(-)}$ и $L_2$.
    Следовательно, $\dim L^{(-)}=\dim L_1=m$. Аналогично показываем, что $\dim L^{(+)} = r-m$. А так как по определению $\dim L^{(-)}$ и $\dim L^{(+)}$ не зависят от базиса, то получаем требуемое.
\end{proof}